\section*{الفصل \ref{ch:b2:two-wave-interference} --- تداخل موجتين}

\begin{solution}{exo:b2:two-wave-interference:1}
$i = \lambda D/a = \qty{2.4}{mm}$؛ و $x/i = 4.2$؛ والهدبة المظلمة العاشرة عند $9.5i = \qty{22.5}{mm}$؛
ومن أجل $a = \qty{2}{mm}$: \qty{0.47}{mm}؛ وفي الماء $\lambda/n$: \qty{1.8}{mm}.
\end{solution}

\begin{solution}{exo:b2:two-wave-interference:2}
$1.25I_0 \pm I_0$: $V = 0.8$. و $1.01I_0 \pm 0.2I_0$: $V = 0.2$ --- إذ يتبع حدّ التداخل
نسبة \emph{السعتين}، $1/10$، لا نسبة الشدتين. وأما المرشّح:
$V = 2\sqrt{0.5}/1.5 = 0.94$.
\end{solution}

\begin{solution}{exo:b2:two-wave-interference:3}
$2ne = \qty{665}{nm}$: مضيء عند $665/(p + \tfrac12)$: \qty{443}{nm} (أزرق بنفسجي)؛
ومظلم عند $665/p$: \qty{665}{nm} (أحمر): فالغشاء أزرق. ومن أجل $e = \qty{1}{\micro m}$: $2ne = \qty{2660}{nm}$،
ويكون مضيئًا عند $760$، و $591$، و $484$، و \qty{409}{nm} --- أي أربعة ألوان معًا، فيُرى
أبيض ورديًّا باهتًا.
\end{solution}

\begin{solution}{exo:b2:two-wave-interference:4}
$r_p = \sqrt{p\lambda R}$: $1.09$، $1.53$، \qty{1.88}{mm}؛ و $p = r^2/\lambda R = 85$ حلقة ضمن
\qty{1}{cm}؛ والمركز مظلم: إذ السماكة معدومة وثمة انعكاس إشارة واحد. وأما الماء:
فتُقسَم أنصاف الأقطار على $\sqrt n$، وتصير $98$ حلقة.
\end{solution}

\begin{solution}{exo:b2:two-wave-interference:5}
(أ) $i = 0.8$ و \qty{1.4}{mm}؛ والأحمر الثاني \qty{2.8}{mm}، والبنفسجي الثالث \qty{2.4}{mm}.
(ب) والرتبة $p$ الحمراء عند $1.4p$ مليمتر تلاقي الرتبة $p + 1$ البنفسجية عند $0.8(p + 1)$: فمن
$p = 2$ تتراكب الرتب. (ج) و $\delta = ax/D = \qty{2.5}{\micro m}$: والغائب $\lambda = \delta/(p +
\tfrac12)$: $714$، $556$، \qty{455}{nm}. (د) فمركز أبيض؛ ثم هدب حرفها
الداخلي أزرق وحرفها الخارجي أحمر؛ وعند الرتبة الثالثة تختلط الألوان
في بياض.
\end{solution}

\begin{solution}{exo:b2:two-wave-interference:6}
(أ) $b < \lambda L/4a = \qty{59}{\micro m}$. (ب) و $\pi ab/\lambda L = 6.7$: فالتباين $|\sin6.7|/6.7
= 0.06$. (ج) ومنبع عرضه سنتيمتر يمتد على مئات الأبعاد الهدبية من
الإزاحة: فلا هدب. (د) و $a < \lambda/\theta = \qty{61}{\micro m}$.
\end{solution}

\begin{solution}{exo:b2:two-wave-interference:7}
(أ) $a = 2h = \qty{0.4}{mm}$: $i = \lambda D/a = \qty{1.25}{mm}$. (ب) وعند $x = 0$ يتساوى المسيران
ومع ذلك تكون الهدبة مظلمة: فالانعكاس يضيف $\pi$ (أي انعكاس
الإشارة عند الورود المماسّ على الوسط الأكثف). (ج) و $\delta = ax/D <
\qty{50}{\micro m}$: أي $100$ هدبة. (د) والبحر هو المرآة، والهوائي على قمة
الجرف هو الشاشة: فكلما صعد المنبع اجتاحت الهدب وأعطت
ارتفاعه.
\end{solution}

\begin{solution}{exo:b2:two-wave-interference:8}
(أ) $2d = 42\lambda$: $d = \qty{12.4}{\micro m}$. (ب) و $100/42 = \qty{2.4}{mm}$. (ج) و $2nd = p\lambda$:
أي $56$ هدبة. (د) وتتحرك الهدب نحو الطرف البعيد وتتباعد كلما نقص $d$؛
وربع هدبة هو $\lambda/8 = \qty{74}{nm}$ على $d$.
\end{solution}

\begin{solution}{exo:b2:two-wave-interference:9}
(أ) الأول وحده (هواء $\to$ زيت، وهو أكثف)؛ وأما زيت $\to$ ماء فإلى قرينة
أدنى: $\delta = 2ne + \lambda/2$، ومضيء عند $2ne = (p + \tfrac12)\lambda$. (ب) و $2ne =
\qty{1160}{nm}$: فالمتعاضد $773$ (على حرف المرئي)، و \qty{464}{nm} (أزرق)؛
والمتلاشي $1160/p$: \qty{580}{nm} (أصفر): فاللون أزرق. (ج) و $2ne = \lambda/2$: $e = \qty{72}{nm}$.
(د) و $\delta = 2ne\cos r$ ينقص مع الزاوية: فتنتقل أطوال الموجة المتعاضدة
نحو الأزرق. (ه) وعندئذ تنعكس الإشارة في الانعكاسين معًا، فلا
$\lambda/2$ زائدة: فمضيء عند $2ne = p\lambda$، ومظلم عند $(p + \tfrac12)\lambda$ --- وهي
طبقة ربع الموجة المانعة للانعكاس.
\end{solution}

\begin{solution}{exo:b2:two-wave-interference:10}
(أ) $(n - 1)e = \qty{10}{\micro m}$، أي $20$ هدبة، نحو الشقّ المغطّى (حيث
يجب أن يقصر المسير الهندسي \qty{10}{\micro m} ليعود $\delta = 0$):
$x = (n - 1)eD/a = \qty{2}{cm}$. (ب) وتقع الهدبة البيضاء هناك؛ وتبدد
الزجاج ($\Delta n \approx 0.01$ عبر المرئي، أي \qty{0.2}{\micro m} من المسير)
دون $\ell_c$: فتبقى مرئية، ملوّنة قليلًا. (ج) وقِس
إزاحة الهدبة البيضاء: $e = (\text{الإزاحة})\,a/(n - 1)D$. (د) و $e(n - 1)\theta^2
(1 + 1/n)/2 = \qty{63}{nm}$ عند \qty{5}{\degree}: أي ثُمن هدبة.
\end{solution}

\begin{solution}{exo:b2:two-wave-interference:11}
(أ) $\theta = \qty{2.3e-7}{rad}$: $a = 1.22\lambda/\theta = \qty{3.1}{m}$. (ب) وحدّه $1.22\lambda/d =
\qty{0.057}{''}$ فوق قطر النجم، والجوّ يشوّش
حتى $1''$. (ج) و $1.22\lambda/100 = \qty{7e-9}{rad} = \qty{1.4}{mas}$. (د) وتحتاج الهدب إلى
$\delta < \ell_c = \lambda^2/\Delta\lambda = \qty{32}{\micro m}$.
\end{solution}

\begin{solution}{exo:b2:two-wave-interference:12}
(أ) $\delta = 2ne\cos r + \lambda/2$؛ و $\cos r = \sqrt{1 - \sin^2i/n^2} \approx 1 - i^2/2n^2$.
(ب) و $2ne/\lambda = 10186.8$: فالمركز، وفيه $p + \tfrac12$ بين $10186.5$ و
$10187.5$، ليس مضيئًا ولا مظلمًا. (ج) و $i_p^2 = 2n^2[1 - (p + \tfrac12)\lambda/
2ne]$ من أجل $p = 10186$، و $10185$، و $10184$: $i = 0.011$، $0.024$، \qty{0.032}{rad}، و $r = fi
= 5.4$، $11.8$، \qty{15.8}{mm}. (د) والزاوية $i$ وحدها تحدد الرتبة: فكل
نقطة منبع تغذّي الحلقة نفسها عند الزاوية نفسها، فتبقى الحلقات، وهي مرئية
في اللانهاية، مع منبع عريض.
\end{solution}

\begin{solution}{pb:b2:two-wave-interference:1}
\textbf{1.} كل مرآة تصوّر $S$ بالتناظر على البعد $d$ من
مستقيم التماس؛ والصورتان تقعان على الدائرة التي نصف قطرها $d$، على
زاوية $2\alpha$ بينهما: $a = 2\alpha d = \qty{2}{mm}$.

\textbf{2.} $S_1$ و $S_2$ صورتان للاهتزاز نفسه: فقفزات
طورهما واحدة.

\textbf{3.} $i = \lambda(d + D)/a = 589 \times 10^{-9} \times 1.7/2 \times 10^{-3} = \qty{0.50}{mm}$: أي عشر
هدب عبر \qty{5}{mm}.

\textbf{4.} $\delta = ax/(d + D) = \qty{2.9}{\micro m}$ عند الحرف، و $p = 5$: أي دون
$\ell_c/\lambda = 34\,000$ بكثير.

\textbf{5.} الإزاحة $bD'/d = 20 \times 10^{-6} \times 1.7/0.2 = \qty{0.17}{mm}$، أي ثلث
$i$: والتباين $|\operatorname{sinc}(\pi ab/\lambda d)| = \operatorname{sinc}(1.07) = 0.82$.

\textbf{6.} $\pi ab/\lambda d = 5.3$: $|\sin5.3|/5.3 = 0.16$.

\textbf{7.} السعتان $\sqrt{0.9}$ و $\sqrt{0.8}$: $V = 2\sqrt{0.72}/1.7 = 0.998$.

\textbf{8.} $(n - 1)e = \qty{5}{\micro m}$: أي $8.5$ هدبة؛ وتُتبَّع بسهولة بضوء
الصوديوم؛ وفي الضوء الأبيض تتحرك الهدبة البيضاء بمقدار $8.5i = \qty{4.2}{mm}$
وتبقى موجودة (إذ يكلّف تبدد الزجاج كسرًا من
الميكرومتر).

\textbf{9.} هدبة مركزية بيضاء، وهدبتان أو ثلاث ملوّنة على
كل جانب، ثم بياض منتظم.

\textbf{10.} هواء $\to$ زيت يعكس الإشارة، وزيت $\to$ ماء (قرينة أدنى) لا يعكسها:
$\delta = 2ne + \lambda/2$؛ ومضيء من أجل $2ne = (p + \tfrac12)\lambda$.

\textbf{11.} $e = \lambda/4n$: الأحمر \qty{112}{nm}، والأخضر \qty{95}{nm}، والبنفسجي \qty{72}{nm}.

\textbf{12.} $2ne = \qty{1740}{nm}$: المتعاضد $1740/(p + \tfrac12)$: $696$ (أحمر)، و $497$
(أزرق أخضر)، و \qty{387}{nm}؛ والمتلاشي $1740/p$: $580$ (أصفر)، و \qty{435}{nm}:
أي خليط أرجواني من الأحمر والأزرق الأخضر.

\textbf{13.} تتغير السماكة على امتداد الطريق؛ وكل لون مضيء
حيث يكون للمقدار $2ne$ القيمة الصحيحة، فتكون الأشرطة خطوط تساوي
السماكة.

\textbf{14.} تنقص الرتب: فتجري الألوان في التتابع
نحو طرف الغشاء الرقيق؛ وعند $e \to 0$ يترك انعكاس الإشارة الوحيد
$\delta = \lambda/2$ من أجل كل لون: فالسواد --- كما في غشاء الصابون في الهواء.

\textbf{15.} بضعة نانومترات: إذ $2ne \ll \lambda$ من أجل الألوان كلها، فيكون
الانعكاسان على تعاكس الطور ويتلاشيان: فالسواد.

\textbf{16.} $\sin r = \sin60^\circ/1.45$، و $r = \qty{37}{\degree}$، و $\cos r = 0.80$: $2ne\cos r
= \qty{1390}{nm}$، ومضيء عند $928$، و $557$، و \qty{398}{nm}: فأحمر السؤال 12
صار أخضر أصفر --- أي نحو الأزرق.

\textbf{17.} حين يتجاوز $2ne$ المقدار $\ell_c \approx \qty{1}{\micro m}$، أي بعد نحو
\qty{0.4}{\micro m} من الزيت، تتراكب الرتب وتخبو الألوان إلى الرمادي.

\textbf{18.} الشعاعان المنعكسان من الغشاء من أجل أي جهة منبع
يلتقيان على الغشاء: فالهدب موضَّعة هناك، ولا يفعل منبع عريض
متعدد الجهات أكثر من إضافة نقوش متطابقة.

\textbf{19.} $\ell_c/\lambda$: $34\,000$ مع الصوديوم، و $1.7 \times 10^7$ مع الليزر.

\textbf{20.} $b < \lambda d/4a$: أي \qty{15}{\micro m} من أجل المرآتين، و \qty{59}{\micro m} من أجل
الشقّين.

\textbf{21.} ارسم نقطتي منبع: فمن أجل كل منهما يتجمع الشعاعان المنعكسان
عند النقطة نفسها من الغشاء، بالمقدار نفسه $\delta = 2ne\cos r$
قرب الورود الناظمي؛ فلا يتحرك نقش الغشاء، بينما يتحرك نقش يونغ.

\textbf{22.} $\lambda/\Delta\lambda = 589/0.02 \approx 30\,000$.

\textbf{23.} $550/30 \approx 18$ هدبة.

\textbf{24.} الضوء المترابط المتبعثر على حبيبات الشاشة يتداخل
بأطوار عشوائية عند العين: وهو الترقيط.

\textbf{25.} المرآتان: \omterm{prop:b2:two-wave-interference:young}{تقسيم جبهة الموجة}، والهدب في كل مكان
(غير موضَّعة)، وعددها يحدّه $\ell_c/\lambda$، ومنبع عريض يمحوها.
وأما الغشاء: \omterm{prop:b2:two-wave-interference:film}{فتقسيم السعة}، والهدب على الغشاء، والألوان
يحدّها $\ell_c \sim \qty{1}{\micro m}$، ومنبع عريض لا يضرّ.
\end{solution}
